\chapter{Tag 15 — Klassen, Vererbung, Encoder als Bibliothek}

\begin{quote}
Funktionen waren das eine Werkzeug — sie kapseln einen Arbeitsschritt.
Klassen sind das andere: sie kapseln Daten und die Funktionen, die
auf diesen Daten arbeiten, in einem einzigen Begriff. Heute lernst
du, was eine Klasse ist und warum sie sich lohnt; danach machen wir
aus deinem Datamatrix-Encoder eine richtige kleine Bibliothek.
\end{quote}

\section*{Lernziele für heute}

Am Ende von Tag 15 kannst du \dots

\begin{itemize}
\item den Unterschied zwischen einer Klasse (Bauplan) und einer
  Instanz (konkretes Objekt) erklären
\item einen Konstruktor (\texttt{\_\_init\_\_}) lesen und schreiben,
  und sagen was \texttt{self} bedeutet
\item Methoden hinzufügen, die mit den Daten der Instanz arbeiten
\item eine Klasse von einer anderen ableiten (\emph{Vererbung}) und
  Methoden gezielt überschreiben
\item den Datamatrix-Encoder als Subklasse einer abstrakten
  \texttt{Barcode}-Klasse aufbauen, und damit Symbolgrößen von
  10$\times$10 bis 26$\times$26 abdecken
\item die offene Schluss-Aufgabe angehen: EAN-13 ins gleiche
  Klassen-Schema integrieren
\end{itemize}

\paragraph{Material} Laptop mit Thonny, dein Code aus den Etappen
10--14, Bleistift und Karopapier für ein paar UML-artige Skizzen
(keine echten UML-Diagramme, nur Kästchen mit Pfeilen).

% =============================================================
\section{Warum reichen Funktionen nicht? \normalfont(\textit{≈ 10 min})}
% =============================================================

In Etappe 14 hast du den Datamatrix-Encoder mit Funktionen aufgebaut.
Das funktionierte — aber spätestens, wenn man zwei Encoder-Varianten
parallel nutzen will, wird es unhandlich:

\begin{itemize}
\item Jede Funktion bekommt \texttt{n=12}, \texttt{anzahl\_pruefbytes=7},
  \texttt{modulgroesse=25} explizit übergeben. Bei drei Aufrufern in
  einem Skript wird das schnell repetitiv.
\item Wer die Codewords bereits berechnet hat, muss sie für die nächste
  Operation in einer eigenen Variable festhalten und wieder weiterreichen.
\item Wenn jemand einen \emph{anderen} Barcode-Typ einbauen will (zum
  Beispiel EAN-13 aus Etappe 10), gibt es keine gemeinsame Schnittstelle
  — jeder Aufrufer muss wissen, welcher Encoder welche Argumente erwartet.
\end{itemize}

Klassen lösen das Problem mit zwei Ideen:

\begin{enumerate}
\item Daten und Funktionen, die zusammengehören, werden in einem
  Begriff zusammengefasst — der Klasse.
\item Verschiedene Klassen, die das Gleiche erfüllen, können von
  einer gemeinsamen Basis erben — und sind dann unter einer
  Schnittstelle austauschbar.
\end{enumerate}

% =============================================================
\section{Klassen-Werkstatt: geometrische Formen \normalfont(\textit{≈ 30 min})}
% =============================================================

Bevor wir Datamatrix anfassen, üben wir Klassen an etwas
Einfacherem. Geometrische Formen eignen sich gut: jede Form hat
einen Namen und eine Fläche, aber die Flächenformel ist je nach
Form anders. Genau die Konstellation, in der Klassen glänzen.

\begin{pythoneinheit}{1}{Eine Basisklasse \texttt{Form}}
\pythoncodeteil{tag15/formen.py}{basisklasse}

\textbf{Was passiert hier?}

\begin{itemize}
\item \texttt{class Form:} eröffnet einen neuen \emph{Bauplan}. Er
  beschreibt, was alle Formen gemeinsam haben.
\item \texttt{\_\_init\_\_(self, name)} ist der \emph{Konstruktor} —
  er wird automatisch aufgerufen, wenn jemand eine Instanz erzeugt:
  \texttt{f = Form("Quadrat")}.
\item \texttt{self} ist die gerade entstehende Instanz. Wenn du in
  einer Methode \texttt{self.name = "Kreis"} schreibst, hängst du
  ein Datenfeld an genau diese Instanz.
\item \texttt{flaeche()} ist eine \emph{abstrakte} Methode: die
  Basisklasse weiß keine Formel, also wirft sie einen Fehler. Jede
  Unterklasse muss \texttt{flaeche()} selbst beisteuern.
\end{itemize}
\end{pythoneinheit}

\begin{pythoneinheit}{2}{Drei Unterklassen}
\pythoncodeteil{tag15/formen.py}{unterklassen}

Drei konkrete Formen — Kreis, Rechteck, Dreieck. Jede schreibt in den
Klammern hinter dem Klassennamen, von wem sie erbt: \texttt{class
Kreis(Form)} bedeutet \emph{Kreis ist eine spezielle Form}.

\texttt{super().\_\_init\_\_("Kreis")} ruft den Konstruktor der
Elternklasse auf — sonst wäre der \texttt{name}-Wert nie gesetzt.
Nach diesem Aufruf hängen wir noch unsere eigene Eigenschaft
\texttt{self.radius} an.

\texttt{flaeche()} überschreibt die Stub-Version aus der Basisklasse
mit der echten Formel. Der Begriff dafür ist \emph{Polymorphismus}
(Vielgestaltigkeit): derselbe Methodenname, aber jede Klasse hat
ihre eigene Implementierung.
\end{pythoneinheit}

\begin{pythoneinheit}{3}{Vererbung in zweiter Stufe — Quadrat}
\pythoncodeteil{tag15/formen.py}{weitervererbung}

Ein Quadrat IST ein spezielles Rechteck. Also erben wir nicht von
\texttt{Form}, sondern von \texttt{Rechteck} — und sparen uns die
\texttt{flaeche()}-Methode komplett, weil sie schon in der
Eltern­klasse stimmt.

Das ist der zentrale Vorteil von Vererbung: gemeinsamer Code wird
nur \emph{einmal} geschrieben. Spezialisierungen passen nur das an,
was sich wirklich unterscheidet.
\end{pythoneinheit}

\begin{pythoneinheit}{4}{Demo: Polymorphismus in einer Liste}
\pythoncodeteil{tag15/formen.py}{demo}

Wir packen vier verschiedene Formen in eine Liste und behandeln sie
alle gleich — die for-Schleife weiß nicht, ob sie gerade einen Kreis
oder ein Dreieck vor sich hat. Jede Instanz weiß selbst, wie sie ihre
Fläche berechnet. Genau das ist Polymorphismus in Aktion.
\end{pythoneinheit}

\begin{bleistiftuebung}{1}{Klassendiagramm zeichnen}
Skizziere auf Karopapier die vier Klassen als Kästchen, mit Pfeilen
nach oben für „erbt von":

\begin{center}
\texttt{Form $\leftarrow$ Kreis} \\
\texttt{Form $\leftarrow$ Rechteck $\leftarrow$ Quadrat} \\
\texttt{Form $\leftarrow$ Dreieck}
\end{center}

In jedes Kästchen kommen die Methoden, die diese Klasse selbst
beisteuert — \texttt{flaeche()} bei den konkreten Klassen,
\texttt{beschreibe()} oben bei \texttt{Form}. So siehst du auf einen
Blick, wer was ererbt und wer was überschreibt.
\end{bleistiftuebung}

\begin{bleistiftuebung}{2}{Eine fünfte Form ergänzen}
Schreibe eine Klasse \texttt{Ellipse(Form)} mit zwei Halbachsen
\texttt{a} und \texttt{b}. Die Flächenformel ist $\pi \cdot a \cdot b$.
Füge eine Instanz davon zur Demo-Liste in \texttt{formen.py} hinzu
und prüfe, dass die Gesamtfläche stimmt.
\end{bleistiftuebung}

% =============================================================
\section{Auf Barcodes übertragen — die abstrakte \texttt{Barcode}-Klasse \normalfont(\textit{≈ 25 min})}
% =============================================================

Jetzt der Sprung zurück zu Datamatrix — mit allem, was du soeben
über Klassen gelernt hast. Die Idee: \emph{ein Barcode ist im
Wesentlichen eine 2D-Matrix aus schwarzen und weißen Modulen}.
Datamatrix, EAN-13, QR-Code — sie alle erfüllen diese
Beschreibung. Also bauen wir eine abstrakte Basisklasse, die genau
das fordert.

\begin{pythoneinheit}{5}{Die abstrakte Basisklasse \texttt{Barcode}}
\pythoncodeteil{tag15/encoder_klasse.py}{barcode}

\texttt{bitmap()} ist die einzige Methode, die jede Subklasse selbst
liefern muss. Den PNG-Export (\texttt{save}) erbt sie geschenkt von
der Basisklasse. Das heißt: wer einen neuen Barcode-Typ einbaut,
muss sich nur um die \emph{Inhalts-Logik} kümmern — der Pixel-Kram
ist schon erledigt.
\end{pythoneinheit}

\begin{pythoneinheit}{6}{\texttt{DatamatrixEncoder} als Subklasse}
\pythoncodeteil{tag15/encoder_klasse.py}{datamatrix}

Wir nehmen die Funktionen aus Etappe 14 und ziehen sie als Methoden
in die Klasse. Die wichtigsten Neuerungen gegenüber Etappe 14:

\begin{itemize}
\item \texttt{SYMBOL\_TABLE}: ein Dictionary mit allen unterstützten
  Symbol-Größen und ihren Daten/ECC-Verhältnissen. Aus ISO 16022.
\item \texttt{\_waehle\_groesse()}: wenn der Aufrufer keine Größe
  angibt, sucht die Klasse die kleinste passende Symbolgröße automatisch.
\item \texttt{\_padden()}: für längere Texte werden Pad-Bytes
  ergänzt. Erstes Pad ist 129; ab dem zweiten kommt ein Pseudo-Zufall
  ins Spiel (sonst würden lange monotone Padding-Felder Scanner
  irritieren).
\item \texttt{bitmap()} ist die einzige Methode, die wir der
  Basisklasse schulden — sie ruft intern \texttt{\_matrix\_bauen()}
  auf, das die Logik aus Etappe 14 in einer Methode bündelt.
\end{itemize}
\end{pythoneinheit}

\begin{pythoneinheit}{7}{Demo: Texte unterschiedlicher Länge}
\pythoncodeteil{tag15/encoder_klasse.py}{demo}

Erwartete Ausgabe:

\begin{minted}{text}
DatamatrixEncoder(text='GRETA', 12×12, 5+7 Codewords)  →  greta.png
DatamatrixEncoder(text='BLAUB', 12×12, 5+7 Codewords)  →  blaub.png
DatamatrixEncoder(text='BLAUBE', 14×14, 8+10 Codewords)  →  blaube.png
DatamatrixEncoder(text='BLAUBEERE', 16×16, 12+12 Codewords)  →  blaubeere.png
DatamatrixEncoder(text='PRAKTIKUMSWOCHE', 18×18, 18+14 Codewords)  →  praktikumswoche.png
\end{minted}

Die Klasse wählt für jeden Text automatisch die kleinste Symbolgröße.
Mit \texttt{pylibdmtx} kannst du alle fünf PNGs prüfen — sie sollten
exakt den Text zurückliefern.
\end{pythoneinheit}

\begin{bleistiftuebung}{3}{Symboltabelle nachvollziehen}
In der Klasse steht \texttt{SYMBOL\_TABLE = \{12: (5, 7), 14: (8, 10),
\dots\}}. Was bedeutet jeder Eintrag? Rechne für mindestens drei
Größen aus, wieviele Module die Innenfläche hat $((n-2)^2)$ und ob
das mit der angegebenen Anzahl Codewords ($\cdot 8$ Bits) hinkommt.
\end{bleistiftuebung}

% =============================================================
\section{Schluss-Aufgabe — EAN-13 ins Schema integrieren \normalfont(\textit{≈ 30 min})}
% =============================================================

Du hast jetzt eine abstrakte \texttt{Barcode}-Klasse und mit
\texttt{DatamatrixEncoder} eine konkrete Implementierung davon.
Die letzte Aufgabe der Etappe ist offen formuliert — und damit
mehr als ein Schreibübung: sie ist der Test, ob du das
Klassen-Konzept wirklich verstanden hast.

\begin{aufgabe}{✏}{Schluss-}{Aufgabe}{cBleistift}{\bfseries
EAN-13 als zweite Subklasse von \texttt{Barcode}.\par}

In Etappe 10 hast du einen vollständigen EAN-13-Encoder mit Funktionen
gebaut: er nimmt 12 Ziffern, berechnet die Prüfziffer, ermittelt für
jede Ziffer das L/G/R-Pattern und setzt am Ende ein 95-Modul breites
Strichmuster zusammen. Die Logik liegt in
\texttt{latex/code/tag10/ean13\_encoder.py}.

Deine Aufgabe: baue daraus eine Klasse \texttt{EAN13Encoder}, die
\emph{genauso} von \texttt{Barcode} erbt wie \texttt{DatamatrixEncoder}.

\paragraph{Was du dafür brauchst}
\begin{itemize}
\item Konstruktor \texttt{\_\_init\_\_(self, ziffern)}, der den Klartext
  (12 oder 13 Ziffern) speichert und ggf. die Prüfziffer ergänzt.
\item Methode \texttt{bitmap()}, die eine 2D-Matrix mit dem 95-Modul-
  Strichmuster liefert. Die Höhe darfst du frei wählen
  (z.\,B. 50 Module hoch — alle Reihen sind identisch, weil
  EAN-13 ein 1D-Code ist und nur in einer Richtung Information trägt).
\item Keine \texttt{save}-Methode — die erbst du von
  \texttt{Barcode}.
\end{itemize}

\paragraph{Was die Aufgabe so spannend macht}

EAN-13 und Datamatrix sehen aus wie unterschiedliche Welten. Trotzdem
sind beide \emph{Barcodes} im Sinne unserer Basisklasse — beide
liefern eine Modul-Matrix. Wenn deine \texttt{EAN13Encoder}-Klasse
fertig ist, kannst du beide Encoder austauschbar verwenden:

\begin{minted}{python}
codes = [
    DatamatrixEncoder("GRETA"),
    EAN13Encoder("400123456789"),
]
for c in codes:
    c.save(f"{c.name}.png")
\end{minted}

Die Schleife weiß nicht, welcher Code-Typ sie gerade verarbeitet.
Sie spricht jeden über die Basisklasse an. Genau das ist die Macht
der Vererbung.

\paragraph{Tipps}
\begin{itemize}
\item Achte auf Kommentare in deiner Klasse — was ist
  \emph{ererbt}, was musst du selbst implementieren?
\item Teste, ob deine PNG-Datei mit einer Smartphone-EAN-Scan-App
  lesbar ist. Tipp aus Etappe 9: Quiet-Zone und Modul-Höhe
  beachten.
\item Wenn du steckenbleibst, schau in die Lösung am Ende des
  Buches — aber versuche es vorher zumindest 30 Minuten allein.
\end{itemize}

\end{aufgabe}

% =============================================================
\section{Reflexion und Brücke zu Tag 16 \normalfont(\textit{≈ 10 min})}
% =============================================================

Du hast jetzt:

\begin{itemize}
\item Klassen, Konstruktoren, \texttt{self}, Methoden verstanden
\item den Unterschied zwischen Klasse (Bauplan) und Instanz
  (konkretes Objekt) verinnerlicht
\item Vererbung in zwei Stufen gesehen
  (\texttt{Form} $\leftarrow$ \texttt{Rechteck} $\leftarrow$
  \texttt{Quadrat})
\item Polymorphismus angewandt: Liste mit verschiedenen Formen,
  alle einheitlich behandelt
\item Datamatrix-Encoding zu einer aufgeräumten Klasse umgebaut,
  die Symbolgrößen 10$\times$10 bis 26$\times$26 abdeckt
\item den Bogen zurück zu Etappe 10 geschlagen — EAN-13 wird zum
  Geschwister von Datamatrix, beide unter dem gemeinsamen Dach
  \texttt{Barcode}
\end{itemize}

\subsection*{Was Klassen jetzt können, was Funktionen nicht
konnten}

\begin{itemize}
\item \textbf{Daten und Verhalten zusammenhalten}: ein
  \texttt{DatamatrixEncoder} weiß seinen Text, seine Größe, seine
  Codewords, seine Matrix. Funktionen mussten das alles als
  Argumente reichen.
\item \textbf{Schnittstellen abstrahieren}: jeder Code, der
  \texttt{Barcode} erwartet, akzeptiert sowohl Datamatrix als auch
  EAN-13. Funktionen hätten dafür eine Wrapper-Funktion gebraucht.
\item \textbf{Spezialisieren ohne Kopieren}: eine Subklasse erbt
  alles aus der Eltern­klasse und passt nur an, was sich
  unterscheidet.
\end{itemize}

\subsection*{Vorschau Tag 16}

Bonus-Sandkasten: Decoder + sichtbare Reed-Solomon-Fehlerkorrektur.
Wir nehmen ein Datamatrix-PNG, übermalen einzelne Module mit Tipp-Ex
oder schwarzem Stift — und gucken zu, wie der Decoder die Fehler
trotzdem repariert. So spürst du den ECC-Anteil zum ersten Mal als
echte Sicherheitsreserve, nicht nur als „diese Bytes sind die ECC".
